% from main.tex
%%%%%%%%%%  scan idx 523  |  folio 510  %%%%%%%%%%
\subsection*{3. Fibrés projectifs et fibrés en drapeaux}

\textbf{Proposition 3.1.} Soient $S$ un schéma noethérien, $E$ un $\mathcal O_S$-Module localement libre de rang $n$, $\pi=(p_1,\ldots,p_k)$ une suite d'entiers tels que $\sum_i p_i=n$, et $X=D_\pi(E)$ le schéma des drapeaux de type $\pi$ de $E$. Alors l'homomorphisme canonique :
\[
K_\bullet(S)\otimes_{K^\circ(S)}K^\circ(X)\longrightarrow K_\bullet(X)
\]
est un isomorphisme.

On reprend les notations de VI 4.5. On note donc $L_i$ les quotients canoniques de l'image inverse de $E$ sur $X$, $F_i$ le noyau de $L_i\to L_{i-1}$ ($i=1,\ldots,k$), et $\lambda_i^{(j)}$ la classe dans $K^\circ(X)$ de $\Lambda^i(F_j)$. Alors les $\lambda_i^{(j)}$ engendrent $K^\circ(X)$ comme $K^\circ(S)$-algèbre (VI 4.6). De même, les images des $\lambda_i^{(j)}$ dans $K^\circ(X_s)$ engendrent $K^\circ(X_s)$ comme algèbre sur $K^\circ(k(s))=\mathbb Z$ pour tout point $s$ de $S$. Comme $X_s$ est un schéma de drapeaux sur un corps, c'est un schéma régulier, donc l'homomorphisme $K^\circ(X_s)\to K_\bullet(X_s)$ est un isomorphisme, si bien que l'homomorphisme $K^\circ(X)\to K_\bullet(X_s)$ est surjectif. D'après 1.3, cela entraîne que l'homomorphisme:
\[
K_\bullet(S)\otimes_{K^\circ(S)}K^\circ(S)\longrightarrow K_\bullet(X)
\]
est surjectif.

Pour montrer l'injectivité, on peut se ramener d'abord au cas du fibré en drapeaux complet $D(E)$. En effet, supposons que
\[
K_\bullet(S)\otimes_{K^\circ(S)}K^\circ(D(E))\longrightarrow K_\bullet(D(E))
\]
soit injectif. Comme $K_\bullet(S)\to K_\bullet(D(E))$ se factorise par $K_\bullet(X)$, cet
